\documentclass[11pt,a4paper]{article}
\usepackage[T1]{fontenc}
\usepackage{lmodern,amsmath,amssymb,amsthm,mathtools,booktabs,microtype}
\usepackage[margin=27mm]{geometry}
\usepackage[colorlinks=true,linkcolor=blue,citecolor=blue,urlcolor=blue]{hyperref}
\setlength{\emergencystretch}{2em}
\newtheorem{theorem}{Theorem}[section]
\newtheorem{lemma}[theorem]{Lemma}
\newtheorem{corollary}[theorem]{Corollary}
\newtheorem{proposition}[theorem]{Proposition}
\newcommand{\Q}{\mathbb Q}\newcommand{\Z}{\mathbb Z}\newcommand{\R}{\mathbb R}
\DeclareMathOperator{\Tr}{Tr}
\title{A positive-rank cubic family for admissible detector traces}
\author{}\date{17 November 2026}
\begin{document}\maketitle
\begin{abstract}
The admissible linear square class \(1+24V\) turns a trace-conic
criterion into a cubic curve with a rational section. We determine
all rational nonsingular weights where this section is torsion, using
exact division-polynomial identities and projective finite-field
certificates. In particular it has infinite order throughout the
normalized physical range \(0<w\le1/4\). Arithmetic progressions of
its multiples produce infinitely many rational traces in \((1,2)\),
arbitrarily close to the reference point at any fixed finite set of
places. These freely varying traces are not identified with prescribed
traces of the original elliptic detectors.
\end{abstract}

\section{The cubic trace family}

For a rational weight \(w\), consider
\begin{equation}\label{eq:trace}
             Y^2=(1+24V)\bigl(V^2+w(2-V)\bigr).
\end{equation}
The changes \(x=1+24V\), \(y=24Y\) give
\begin{equation}\label{eq:E}
 E_w:\quad y^2=x^3-(2+24w)x^2+(1+1176w)x.
\end{equation}
The points
\[
                        P=(49,336),\qquad T=(0,0)
\]
lie on every member of the family. Its discriminant is
\begin{equation}\label{eq:disc}
                 \Delta_w=9216w(w-8)(1+1176w)^2.
\end{equation}
For nonsingular weights, \(T\) has order two.

\begin{theorem}\label{thm:rank}
For rational \(w\notin\{0,8,-1/1176\}\), the point \(P\) is
torsion if and only if
\[
 w\in\left\{\frac{296}{49},-\frac{96}{49},\frac{100}{49}\right\}.
\]
Its orders at these three weights are respectively three, three and
four. Consequently, for every rational \(0<w\le1/4\), the smooth
curve \(E_w\) has positive Mordell--Weil rank over \(\Q\).
\end{theorem}

\section{Exact torsion exclusion}

We give the algebraic specification and finite certificates used to
prove Theorem~\ref{thm:rank}. Let
\[
 a_2=-2-24w,\quad a_4=1+1176w,\quad
 b_2=4a_2,\quad b_4=2a_4,\quad b_6=0,\quad b_8=-a_4^2.
\]
Let \(\psi_m=\psi_m(P)\) denote the standard division polynomial
evaluated at \(x=49,y=336\). The initial values and recurrence are
\begin{align*}
 \psi_0&=0,&\psi_1&=1,&\psi_2&=672,\\
 \psi_3&=3x^4+b_2x^3+3b_4x^2+b_8,\\
 \psi_4&=2y\bigl(2x^6+b_2x^5+5b_4x^4
                       +10b_8x^2+b_2b_8x+b_4b_8\bigr),\\
 \psi_{2m+1}&=\psi_{m+2}\psi_m^3-\psi_{m-1}\psi_{m+1}^3,\\
 \psi_{2m}&=\frac{\psi_m}{672}
           \bigl(\psi_{m+2}\psi_{m-1}^2
                         -\psi_{m-2}\psi_{m+1}^2\bigr).
\end{align*}
Every resulting expression used below is an integral polynomial in
\(w\). The standard vanishing criterion is
\(\psi_m(P)=0\) if and only if \([m]P=O\) on a nonsingular
fiber; see \cite[Chapter III, Exercise 3.7]{Silverman}.

Put
\[
 L_8=49w-100,\qquad L_{10}=1,\qquad
 L_{12}=(49w-100)(49w-296)(49w+96).
\]
For \(m\in\{8,10,12\}\), define \(Q_m\) to be the primitive
integer polynomial with positive leading coefficient obtained from
\(\psi_m/L_m\). The exact divisions have zero remainder.
The following table gives projective no-root certificates:
\begin{center}
\begin{tabular}{rrrrr}\toprule
\(m\)&\(\deg Q_m\)&prime \(p\)&leading coefficient modulo \(p\)
  &roots in \(\mathbb P^1(\mathbb F_p)\)\\\midrule
8&14&5&1&0\\
10&24&877&322&0\\
12&32&449&179&0\\\bottomrule
\end{tabular}
\end{center}

The companion certificate \cite{Certificate} gives the full integer
coefficient arrays, exact division identities, and all affine residue
values. Its verifier reconstructs the displayed recurrences with
standard-library rational polynomial arithmetic, checks the quotient
identities, and evaluates at every residue. It also checks that the
three moduli are prime and that the leading coefficients are nonzero.
Thus infinity is tested explicitly, rather than being omitted by an
affine search.

As a second finite verification, the same certificate records an exact
factorization of each \(Q_m\). Every nonlinear factor has a projective
no-root certificate at one of the smaller primes \(5,11,13\).
There are eleven distinct nonlinear factors, of degrees
\[
                 2,4,8,3,3,6,12,2,4,8,16.
\]
The verifier checks all factor products exactly and their leading
coefficients as well; it uses no factorization routine. An additional checker computes $[8]P,[10]P,[12]P$ directly by
tangent and chord addition over $\Q(w)$. Their reduced abscissa
denominators are nonzero rational constant multiples of
$\psi_8^2,\psi_{10}^2,\psi_{12}^2$, respectively. This independently
checks the multiplication denominators without using the recurrence.
Neither certificate relies on floating-point roots, a rank database or an
unverified Mordell--Weil basis.

\begin{lemma}\label{lem:projective}
Let \(Q\in\Z[w]\) have leading coefficient nonzero modulo a prime
\(p\). If its reduction has no root in \(\mathbb F_p\), then
\(Q\) has no rational root.
\end{lemma}
\begin{proof}
Write a proposed root as \(a/b\) in lowest terms. If \(p\nmid b\),
reduction gives an affine root. If \(p\mid b\), then \(p\nmid a\),
and the homogeneous equation of degree \(\deg Q\) reduces to the
nonzero leading coefficient times \(a^{\deg Q}\). This is impossible.
\end{proof}

\begin{proof}[Proof of Theorem~\ref{thm:rank}]
Since \(E_w\) contains a rational point of order two, Mazur's torsion
classification \cite[Theorem 8; III, Theorem 5.1]{Mazur} implies that the order of any
rational torsion point divides one of \(8,10,12\). Indeed the only
additional possible orders in the general classification are seven
and nine; together with the rational two-torsion they would give
torsion of order fourteen or eighteen, which the classification excludes.
The point \(P\) is neither the identity nor two-torsion, because its
ordinate is \(336\ne0\).

The exact quotient certificates and Lemma~\ref{lem:projective} show
that the only rational weights at which any of
\(\psi_8,\psi_{10},\psi_{12}\) vanishes are the three listed roots
of \(L_8L_{12}\). At the first two, \([2]P=-P\); at the third,
\([2]P=T\). These identities follow directly from the group law,
or from the exact formulas
\[
 \psi_3=-576(49w-296)(49w+96),\qquad
                         x([2]P)=\frac{(49w-100)^2}{784}.
\]
All three fibers are nonsingular. Their orders are therefore exactly
three, three and four. They are outside \(0<w\le1/4\), and the
singular weights are outside that interval too. This proves the theorem.
\end{proof}

\section{Approximation and infinitely many free trace values}

\begin{corollary}\label{cor:density}
Fix a rational \(0<w\le1/4\), a finite set of rational primes, and
arbitrarily small neighborhoods of \(P\) in the corresponding
\(E_w(\Q_p)\). There are infinitely many rational points on
\eqref{eq:trace} with \(1<V<2\), whose points on \(E_w\) lie in
all those neighborhoods and are arbitrarily close to \(P\) over
\(\R\). A finite set of nonzero rational-function guards on
\(E_w\) can simultaneously be avoided.
\end{corollary}
\begin{proof}
The discriminant of the quadratic factor of the cubic in
\eqref{eq:E} is \(576w(w-8)<0\). Thus the cubic has one real root
and \(E_w(\R)\) is connected, topologically a circle.
By Theorem~\ref{thm:rank}, \(P\) has infinite order, so the closure
of its cyclic subgroup is the whole real group: an infinite closed
subgroup of a circle is the circle itself.

For each chosen finite place, take an open subgroup \(U_p\) such
that \(P+U_p\) lies inside the prescribed neighborhood. The compact
group \(E_w(\Q_p)\) has finite quotient by \(U_p\). Choose a
positive integer \(M\) annihilating the class of \(P\) in all
these finite quotients. Every point \([1+Mk]P\) then satisfies all
finite-place requirements. The point \([M]P\) still has infinite
order, so these points remain dense in the real group.

The abscissa map is unramified at \(P\), because its ordinate is
nonzero. The real neighborhoods of \(P\) therefore contain an open
arc with \(25<x<49\), equivalently \(1<V<2\).
Infinitely many points of the progression lie in that arc, arbitrarily
close to \(P\). Every nonzero rational function on a smooth curve
has finitely many zeros and poles. Removing the finitely many points
excluded by the guards leaves infinitely many suitable points.
\end{proof}

The trace map \(V=(x-1)/24\) has finite fibers, so the corollary
also gives infinitely many distinct rational values of \(V\).
The real and finite approximation statement uses only one fixed
weight at a time; independent copies give independently chosen traces
for any fixed finite list of weights.

\section{Scope for the defining-polynomial construction}

For a positive normalized general input \(u\), the physical weight
\(w=u/(1+u)^2\) lies in \((0,1/4]\), so the entire family above
is available. The conditional preparation can likewise choose small
positive weights. The admissible class \(1+24V\) has no bad-pole
height loss in the trace criterion \cite{Trace}.

The original detector, however, prescribes
\(V=\Tr(W^{-1})\) as a rational function of its elliptic point tuple
and common slope. Corollary~\ref{cor:density} does not identify its
freely chosen trace with that value. It does not prove simultaneous
matching of two actual traces, a new defining-polynomial count, or a
rational section over the full source space. These remain separate
arithmetic requirements.

\begin{thebibliography}{9}\raggedright
\bibitem{Mazur}B.~Mazur, \emph{Modular curves and the Eisenstein ideal},
Publications Math\'ematiques de l'IH\'ES \textbf{47} (1977), 33--186,
Theorem 8 and III, Theorem 5.1.
\url{https://www.numdam.org/item/PMIHES_1977__47__33_0.pdf}.
\bibitem{Silverman}J.~H.~Silverman, \emph{The Arithmetic of Elliptic
Curves}, second edition, Graduate Texts in Mathematics 106, Springer,
2009, Chapter III, Exercise 3.7.
\bibitem{Certificate}Exact arithmetic certificate and independent
standard-library verifier,
\nolinkurl{computation/pole_parameters/cubic_trace_rank_v57/division_data.json},
\nolinkurl{computation/pole_parameters/cubic_trace_rank_v57/check.py}.
\bibitem{Trace}\emph{A rational conic criterion using the actual
detector trace},
\nolinkurl{trace_conic_criterion_20261117_v57.tex}.
\end{thebibliography}
\end{document}
